\chapter[PROJETO DO SISTEMA]{Projeto do Sistema}
\label{cap:projeto}

Este capítulo apresenta o projeto do sistema desenvolvido, denominado WebGate.
São descritos os requisitos que orientaram o desenvolvimento, a arquitetura e
seus componentes, a sequência de etapas do \textit{pipeline}, o modelo de dados
e as principais decisões de projeto. As heurísticas de classificação usadas nos
filtros foram apresentadas na Seção~\ref{sec:heuristicas}, e a forma como o
projeto foi codificado é descrita no Capítulo~\ref{cap:implementacao}.

% -----------------------------------------------------------------------------
\section{Requisitos}
\label{sec:requisitos}

O objetivo do sistema é oferecer a empresas incubadas uma lista de editais de
fomento com inscrições abertas, encontrada automaticamente na Web e atualizada
sem intervenção do usuário. A partir desse objetivo, foram definidos os
requisitos funcionais da Tabela~\ref{tab:rf} e os requisitos não funcionais da
Tabela~\ref{tab:rnf}.

\begin{table}[htbp]
\centering
\caption{Requisitos funcionais}
\label{tab:rf}
\small
\begin{tabularx}{\textwidth}{l X}
\toprule
Código & Descrição \\
\midrule
RF01 & Gerar consultas de busca a partir de termos do domínio dos editais de
       fomento. \\
RF02 & Recuperar páginas candidatas por meio de uma API de busca na Web. \\
RF03 & Baixar o conteúdo de cada página candidata, em HTML ou PDF, e extrair
       seu texto. \\
RF04 & Descartar páginas agregadoras, que listam vários editais em vez de
       descrever um. \\
RF05 & Descartar páginas cujo conteúdo não corresponde ao escopo do sistema. \\
RF06 & Identificar o prazo de inscrição e descartar editais sem prazo ou com
       prazo vencido. \\
RF07 & Descartar páginas duplicadas, por URL e por conteúdo. \\
RF08 & Armazenar os editais aceitos em um banco de dados. \\
RF09 & Atualizar o banco de dados automaticamente, em intervalos regulares. \\
RF10 & Exibir ao usuário, em uma página web, os editais com inscrições abertas,
       com título, \textit{link}, domínio e prazo. \\
RF11 & Permitir que o usuário solicite a lista mais recente de editais por meio
       de um botão ``Atualizar''. \\
RF12 & Registrar a decisão tomada sobre cada página, aceita ou não, e o motivo
       da decisão. \\
\bottomrule
\end{tabularx}
\fonte{Elaborado pelo autor.}
\end{table}

\begin{table}[htbp]
\centering
\caption{Requisitos não funcionais}
\label{tab:rnf}
\small
\begin{tabularx}{\textwidth}{l X}
\toprule
Código & Descrição \\
\midrule
RNF01 & \textbf{Reprodutibilidade}: toda coleta deve poder ser reprocessada
        offline, com o mesmo resultado, independentemente de mudanças na Web. \\
RNF02 & \textbf{Configurabilidade}: pesos e limiares das heurísticas devem ser
        alteráveis sem modificação do código das heurísticas. \\
RNF03 & \textbf{Interpretabilidade}: cada decisão do sistema deve poder ser
        justificada pelas regras que a produziram. \\
RNF04 & \textbf{Economia de requisições}: reabrir a página não deve gerar
        requisição ao servidor, e cada coleta deve usar um número limitado de
        chamadas à API de busca. \\
RNF05 & \textbf{Responsividade}: a interface deve ser utilizável em telas de
        celulares, \textit{tablets} e computadores. \\
RNF06 & \textbf{Robustez}: falhas de rede, da API de busca ou do banco de
        dados não devem interromper o servidor. \\
RNF07 & \textbf{Segurança}: conteúdos vindos da Web não devem ser executados na
        página do usuário, e credenciais não devem constar do código-fonte. \\
\bottomrule
\end{tabularx}
\fonte{Elaborado pelo autor.}
\end{table}

% -----------------------------------------------------------------------------
\section{Arquitetura do sistema}
\label{sec:arquitetura}

O sistema implementa a arquitetura em etapas apresentada no Referencial
Teórico, na qual a recuperação de documentos candidatos é delegada a uma API de
busca e a avaliação da relevância é feita posteriormente, por filtros próprios.
As tecnologias empregadas estão descritas na Seção~\ref{sec:tecnologias}. O
sistema é composto por quatro componentes, que se comunicam em um único
sentido, da coleta até o usuário:

\begin{enumerate}
  \item \textbf{\textit{Pipeline} de editais}: executa a coleta e o
    processamento descritos nas Seções~\ref{subsec:fase_coleta}
    e~\ref{subsec:fase_processamento}. Recebe como entrada as respostas da API
    de busca e as páginas baixadas, e produz como saída o \textit{snapshot} da
    coleta, o arquivo de predições e os editais aceitos.
  \item \textbf{Banco de dados}: armazena os editais aceitos por todas as
    coletas, de forma acumulada. É a única fonte de dados da interface.
  \item \textbf{Servidor web}: agenda a execução periódica do
    \textit{pipeline} e expõe os editais do banco por meio de um
    \textit{endpoint} de leitura.
  \item \textbf{Cliente web}: página carregada no navegador do usuário, que
    solicita a lista de editais ao servidor, guarda-a localmente e a apresenta
    na forma de cards.
\end{enumerate}

O fluxo de dados pode ser representado por:

\begin{equation*}
\text{API de busca}
\rightarrow
\textit{Pipeline}
\rightarrow
\text{Banco de dados}
\rightarrow
\text{Servidor web}
\rightarrow
\text{Navegador}
\end{equation*}

Os componentes funcionam em ritmos independentes. O \textit{pipeline} é
disparado pelo servidor a cada 24 horas e grava no banco. O cliente consulta o
servidor somente quando o usuário pressiona o botão ``Atualizar'', e o
servidor responde com o conteúdo do banco naquele momento, sem executar o
\textit{pipeline}. Assim, o tempo de resposta da página não depende da
duração de uma coleta, que envolve dezenas de requisições externas.

Além desses componentes, o sistema inclui um módulo de avaliação, que não faz
parte do fluxo de produção. Ele lê os \textit{snapshots}, as predições e os
rótulos manuais e calcula as métricas utilizadas no procedimento de calibração
(Seção~\ref{sec:calibracao}).

O \textit{pipeline} é dividido em duas fases independentes: coleta e
processamento. A coleta consulta a API de busca, baixa as páginas encontradas e
grava o resultado em disco. O processamento lê esse registro e aplica, em
cadeia, os filtros que decidem se cada página é um edital aberto. A separação
permite reprocessar a mesma coleta quantas vezes forem necessárias, com
configurações diferentes, sem depender do estado da Web no momento do
experimento.

\subsection{Fase de coleta}
\label{subsec:fase_coleta}

A fase de coleta compreende as seguintes etapas:

\begin{enumerate}
  \item \textbf{Geração de consultas}: combinação de termos-base (``edital'',
    ``chamada pública'', ``bolsa pesquisa'', ``financiamento inovação'',
    ``inscrições abertas'', ``chamada aberta'') com o operador \texttt{site:},
    restrito a domínios institucionais (\texttt{gov.br}, \texttt{fapesp.br},
    \texttt{cnpq.br}, \texttt{capes.gov.br}), e com o ano de referência, como
    em \texttt{edital site:fapesp.br 2026}.
  \item \textbf{Priorização das consultas}: as consultas restritas a um
    domínio, com o ano e com o termo ``edital'' recebem prioridade, somada ao
    histórico de desempenho de cada consulta em execuções anteriores.
  \item \textbf{Busca}: cada consulta é enviada à API Tavily \cite{tavily},
    com até 5 resultados por consulta, até o limite de 50 resultados brutos.
  \item \textbf{Deduplicação por URL}: as URLs são normalizadas (sem esquema,
    sem ``\texttt{www.}'', sem porta, sem barra final e sem fragmento) e as
    repetidas são descartadas.
  \item \textbf{Download e registro}: o conteúdo bruto de cada página (HTML ou
    PDF) é baixado com a biblioteca Requests \cite{requests} e gravado no
    \textit{snapshot} da coleta (Seção~\ref{sec:coleta}).
\end{enumerate}

\subsection{Fase de processamento}
\label{subsec:fase_processamento}

A fase de processamento aplica a cada item, nesta ordem:

\begin{enumerate}
  \item \textbf{Extração de texto}: o HTML é convertido em texto com a
    biblioteca Beautiful Soup \cite{beautifulsoup}, com remoção de
    \textit{scripts} e estilos; o PDF tem o texto extraído página a página com a
    biblioteca pypdf \cite{pypdf}. Quando o texto extraído é menor que o resumo
    (\textit{snippet}) devolvido pela API, usa-se o resumo.
  \item \textbf{Filtro de listagem}: descarta páginas agregadoras, que listam
    vários editais em vez de descrever um (Seção~\ref{subsec:listagem}).
  \item \textbf{Filtro de relevância}: descarta páginas cujo vocabulário não
    corresponde ao escopo do sistema (Seção~\ref{subsec:relevancia}).
  \item \textbf{Extração de prazo}: identifica a data-limite de inscrição e
    descarta os itens sem prazo ou com prazo vencido na data da coleta
    (Seção~\ref{subsec:prazo}).
  \item \textbf{Deduplicação por conteúdo}: descarta itens cuja assinatura MD5
    \cite{rivest1992} dos primeiros 1.000 caracteres normalizados já foi vista,
    ou cuja similaridade de Jaccard \cite{jaccard1901,broder1997} entre os
    conjuntos de palavras dos primeiros 1.000 caracteres supera 0,8 em relação
    a um item já aceito.
  \item \textbf{Persistência}: os itens aceitos são gravados em banco de dados
    PostgreSQL \cite{postgresql}, executado em contêiner Docker \cite{docker}.
    Todos os itens, aceitos ou não, são registrados em um arquivo de predições
    com o motivo da decisão.
\end{enumerate}

O desenho em cadeia implica que cada filtro só recebe os itens aprovados pelo
anterior. Um edital descartado indevidamente pelo filtro de listagem, por
exemplo, nunca chega à extração de prazo. Essa propriedade orientou a escolha
das métricas (Seção~\ref{sec:metricas}). Nenhuma etapa decide explicitamente se
a página é um edital: essa função é exercida, de modo indireto, pelos filtros
de listagem e de relevância.

% -----------------------------------------------------------------------------
\section{Modelo de dados}
\label{sec:modelo_dados}

O sistema mantém dados em dois meios: o banco de dados relacional, consultado
pela interface, e arquivos em disco, que registram cada coleta para fins de
reprodutibilidade e avaliação.

\subsection{Banco de dados}

O banco de dados contém uma única entidade, \texttt{Edital}, que representa um
edital aceito pelo \textit{pipeline}. Seus atributos estão na
Tabela~\ref{tab:entidade_edital}. A URL é única, o que impede que o mesmo
edital seja gravado duas vezes por coletas diferentes.

\begin{table}[htbp]
\centering
\caption{Atributos da entidade \texttt{Edital}}
\label{tab:entidade_edital}
\small
\begin{tabularx}{\textwidth}{>{\ttfamily}l l X}
\toprule
\normalfont Atributo & Tipo & Descrição \\
\midrule
id         & inteiro        & Identificador gerado pelo banco (chave primária). \\
title      & texto          & Título da página, devolvido pela API de busca. \\
url        & texto          & Endereço da página. Indexado e único. \\
content    & texto          & Texto extraído da página, usado nas heurísticas. \\
deadline   & data           & Prazo de inscrição estimado pela extração de prazo. \\
domain     & texto          & Domínio da URL, exibido no card. \\
score      & inteiro        & Pontuação de relevância $R$ (Equação~\ref{eq:relevancia}). \\
phrase     & texto          & Frase da página de onde o prazo foi extraído. \\
created\_at & data e hora   & Instante de gravação, em UTC. Corresponde ao dia da coleta. \\
\bottomrule
\end{tabularx}
\fonte{Elaborado pelo autor.}
\end{table}

\subsection{Arquivos de coleta}

Cada coleta gera um diretório próprio, nomeado pelo instante da coleta, com o
conteúdo apresentado na Seção~\ref{sec:coleta}. A
Tabela~\ref{tab:arquivos_coleta} resume os arquivos desse diretório. Os três
primeiros são produzidos pelo \textit{pipeline}; os demais, pelo módulo de
avaliação.

\begin{table}[htbp]
\centering
\caption{Arquivos do diretório de uma coleta}
\label{tab:arquivos_coleta}
\small
\begin{tabularx}{\textwidth}{>{\ttfamily}l X}
\toprule
\normalfont Arquivo & Conteúdo \\
\midrule
raw/            & Conteúdo bruto de cada página, nomeado pelo \textit{hash}
                  SHA-256 do próprio conteúdo. \\
manifest.json   & Instante da coleta, \textit{commit}, versões das
                  bibliotecas, consultas executadas, configuração vigente e
                  metadados de cada item. \\
predictions.csv & Uma linha por item, com as decisões de cada etapa e o motivo
                  da decisão final. \\
labels.csv      & Rótulos manuais de cada item (Tabela~\ref{tab:rotulos}). \\
metrics.csv     & Métricas de classificação de cada etapa. \\
features.csv    & Pontuações das heurísticas para cada item, calculadas sem a
                  cadeia de filtros. \\
\bottomrule
\end{tabularx}
\fonte{Elaborado pelo autor.}
\end{table}

Há ainda um arquivo de estatísticas das consultas, compartilhado entre as
coletas, que acumula para cada consulta o número de usos e uma pontuação de
desempenho. Essa pontuação é usada na priorização das consultas das coletas
seguintes.

% -----------------------------------------------------------------------------
\section{Decisões de projeto}
\label{sec:decisoes}

As decisões a seguir foram tomadas para atender aos requisitos não funcionais.

\textbf{Separação entre coleta e processamento} (RNF01). A coleta e o
processamento são fases distintas, ligadas apenas pelo \textit{snapshot}. O
processamento não acessa a Web: lê o conteúdo gravado e usa como data de
referência o instante registrado no manifesto. Com isso, qualquer coleta
anterior pode ser reprocessada com outra configuração, e o resultado depende
apenas do código e da configuração, e não do dia em que é executado. Essa
propriedade é a base do procedimento de calibração e do teste de regressão.

\textbf{Persistência opcional} (RNF01). O \textit{pipeline} pode ser executado
com ou sem gravação no banco. No modo de reprocessamento, usado pelo módulo de
avaliação, o banco não é acessado, de modo que os experimentos não dependem do
estado do banco nem o alteram.

\textbf{Configuração centralizada e imutável} (RNF02 e RNF03). Todos os pesos,
limiares, vocabulários e opções das heurísticas ficam em uma única estrutura de
configuração, dividida em grupos (listagem, relevância, prazo e deduplicação).
A estrutura é imutável: uma configuração alternativa é criada como cópia da
padrão, com os valores alterados, sem modificar a original. Isso permite que a
análise de sensibilidade e a busca em grade avaliem centenas de configurações
no mesmo processo, sem interferência entre elas. A configuração vigente é
gravada no manifesto de cada coleta.

\textbf{Variantes como opções desligadas por padrão} (RNF01). Toda mudança
estrutural nas heurísticas foi introduzida como uma opção da configuração,
inicialmente desligada. Assim, o comportamento padrão só muda quando a opção é
ligada de forma deliberada, e o teste de regressão
(Seção~\ref{sec:coleta}) pode confirmar que a introdução da opção não alterou
nenhuma decisão.

\textbf{Registro de todas as decisões} (RF12 e RNF03). O arquivo de predições
registra todos os itens, inclusive os descartados, com o motivo do descarte.
Isso permite comparar as decisões do sistema com os rótulos manuais e
identificar em qual etapa cada erro ocorreu.

\textbf{Atualização no servidor e cache no cliente} (RF09 e RNF04). A coleta é
de responsabilidade exclusiva do servidor, que a executa a cada 24 horas. O
cliente guarda no navegador a última lista recebida e só faz nova requisição
quando o usuário pressiona ``Atualizar''. Desse modo, o número de chamadas à
API de busca não depende do número de usuários nem de quantas vezes a página é
aberta.

\textbf{Filtros de exibição separados do \textit{pipeline}}. Ao entregar a
lista à interface, o servidor descarta prazos implausíveis, situados a mais de
cinco anos da data da coleta, e cards repetidos, com o mesmo título e o mesmo
prazo. Esses filtros atuam apenas sobre a exibição e não alteram o banco nem o
\textit{pipeline}. Por isso, não afetam as métricas apresentadas no
Capítulo~\ref{cap:resultados}, que se referem à saída do \textit{pipeline}.
A decisão de mantê-los fora do \textit{pipeline} preservou a comparabilidade
com os resultados da calibração, já congelada quando os filtros foram
introduzidos.

% -----------------------------------------------------------------------------
\section{Interface do usuário}
\label{sec:proj_interface}

A interface é uma página única, de consulta, sem cadastro de usuários e sem
controles de execução. Ela é composta por um cabeçalho, com o nome do sistema
e a frase ``Encontre editais que podem fazer sua ideia ganhar vida'', e por um
quadro de editais, com o título ``Editais encontrados'', o botão
``Atualizar'' e uma grade de cards. Cada card apresenta o título do edital,
com \textit{link} para a página original, o prazo de inscrição e o domínio.
Prazos a 30 dias ou menos do encerramento recebem destaque visual.

O comportamento da página segue três situações:

\begin{enumerate}
  \item \textbf{Primeira visita}: não há lista guardada no navegador. O quadro
    aparece vazio, com uma mensagem que orienta o usuário a pressionar
    ``Atualizar'', e nenhuma requisição é feita.
  \item \textbf{Atualização}: ao pressionar ``Atualizar'', a página solicita a
    lista ao servidor, guarda-a no navegador e redesenha o quadro. Uma linha de
    \textit{status} informa quando o servidor fez a última coleta, quando a
    lista foi carregada e quando está prevista a próxima coleta.
  \item \textbf{Visitas seguintes}: a página exibe a lista guardada, sem
    requisição. Editais cujo prazo venceu depois que a lista foi guardada são
    omitidos.
\end{enumerate}

Para atender ao requisito de responsividade (RNF05), o número de colunas da
grade de cards varia com a largura da tela: uma coluna em celulares, duas a
partir de 576 \textit{pixels}, o que inclui os \textit{tablets}, três a partir
de 992 \textit{pixels} e quatro a partir de 1.400 \textit{pixels}.
